\elementopretextual{Abstract}{abstract}
\begingroup
\singlespacing
\noindent\textbf{Title:} \textit{In silico} screening of compounds selected for the prospection of \textit{Tetraselmis chuii} against protein targets associated with skin aging.\par
\vspace{0.5\baselineskip}
\noindent
Skin aging involves changes in the extracellular matrix and pigmentation, prompting the study of compounds that may interact with proteins associated with these processes. This study applied molecular docking to compare 17 commercial and reference ligands with 17 ligands in the \textit{T. chuii} SPE prospecting group. The 34 ligands were evaluated against six human protein structures: three matrix metalloproteinases (MMP-9, 1GKC; MMP-1, 1HFC; and MMP-2, 8H78), neutrophil elastase (1H1B), hyaluronidase 1 (2PE4), and tyrosinase-related protein 1 (TYRP1, 5M8P). Literature screening considered reports on biomass, cell walls, exopolysaccharides, and extracellular vesicles to distinguish the sources associated with the selected compounds. AutoDock Vina with its Vina scoring function was used for elastase and hyaluronidase; the other four targets were evaluated with AutoDock4 (AD4) scoring and AutoDock4Zn (AD4Zn) maps. The analysis included 204 FEB (free energy of binding) scores, one for each receptor--ligand pair, as well as PyMOL visualization, PLIP-identified contacts, and conformational comparisons with available crystallographic references. Within the SPE prospecting group, glucose-6-sulfate ranked highest for MMP-9 and MMP-1, galactose-6-sulfate for MMP-2, Dha (3-deoxy-D-lyxo-hept-2-ulopyranosaric acid) for elastase and hyaluronidase, and ferulic acid for TYRP1. The leading commercial ligands yielded more negative scores for all six targets. Together, these findings organize ligands by score and contact profiles and guide the selection of SPE group compounds for further investigation. Chemical characterization of the SPE from \textit{T. chuii}, followed by biochemical and cellular assays, will help identify its constituents and assess their potential for skin applications.

\vspace{\baselineskip}
\noindent\textbf{Key words:} molecular docking; \textit{Tetraselmis chuii}; extracellular polymeric substances; skin aging; computational screening.
\endgroup
